% GENERATED FILE. Edit canonical lessons, then run npm run generate:books.
% canonical-source-hash: fnv1a64:7aef5124637ec5d0
% canonical-lessons: TE-C251-manasu, TE-C251-ayasam, TE-C251-kandaram, TE-C251-injeksanu, TE-C251-tika

\chapter{Mind, Breathlessness, Muscle, Injection, Vaccination}
\label{ch:te-mind-breathlessness-muscle-injection-vaccination}
\hlchaptermodality{251}

\begin{chapteropening}
\textbf{By the end of this chapter:} \emph{I can say the mind, breathlessness, a muscle, an injection and a vaccination.}

Complete the last lesson of chapter 251: I can say the mind, breathlessness, a muscle, an injection and a vaccination.
\end{chapteropening}

\section[\texorpdfstring{\te{మనసు} (manasu)}{మనసు (manasu)}]{\te{మనసు} (manasu) — the mind, one's inner feelings}

\label{lesson:TE-C251-manasu}



\begin{quote}
Before the new one: say the Telugu for the thigh, then the Telugu for the waist.
\end{quote}

\subsection*{You'll want to know: \te{మనసు}}
\textbf{\te{మనసు}} — \emph{manasu} — \textquotedblleft{}the mind, one's inner feelings\textquotedblright{}.

\begin{grammarlens}[title={What you've built}]
One more word for this chapter.
\end{grammarlens}

\subsection*{Your turn}
\begin{itemize}
\raggedright
  \item \emph{Say it:} \emph{manasu}
  \item \emph{Say it:} \emph{manasu}, once more
  \item \emph{Say it:} say \emph{manasu}
  \item \emph{Recall:} say the Telugu for the thigh, then the Telugu for the waist, then say \emph{manasu} again
\end{itemize}

\begin{tcolorbox}[breakable,colback=teal!4,colframe=teal!35!black,title={Before you move on}]
What does \te{మనసు} mean? (\textquotedblleft{}The mind, one's inner feelings\textquotedblright{}.) Say it once more.
\end{tcolorbox}
\section[\texorpdfstring{\te{ఆయాసం} (āyāsaṁ)}{ఆయాసం (āyāsaṁ)}]{\te{ఆయాసం} (āyāsaṁ) — breathlessness}

\label{lesson:TE-C251-ayasam}



\begin{quote}
Before the new one: say the Telugu for the waist, then the Telugu for the mind.
\end{quote}

\subsection*{You'll want to know: \te{ఆయాసం}}
\textbf{\te{ఆయాసం}} — \emph{āyāsaṁ} — \textquotedblleft{}breathlessness\textquotedblright{}.

\begin{grammarlens}[title={What you've built}]
One more word for this chapter.
\end{grammarlens}

\subsection*{Your turn}
\begin{itemize}
\raggedright
  \item \emph{Say it:} \emph{āyāsaṁ}
  \item \emph{Say it:} \emph{āyāsaṁ}, once more
  \item \emph{Say it:} say \emph{āyāsaṁ}
  \item \emph{Recall:} say the Telugu for the waist, then the Telugu for the mind, then say \emph{āyāsaṁ} again
\end{itemize}

\begin{tcolorbox}[breakable,colback=teal!4,colframe=teal!35!black,title={Before you move on}]
What does \te{ఆయాసం} mean? (\textquotedblleft{}Breathlessness\textquotedblright{}.) Say it once more.
\end{tcolorbox}
\section[\texorpdfstring{\te{కండరం} (kaṇḍaraṁ)}{కండరం (kaṇḍaraṁ)}]{\te{కండరం} (kaṇḍaraṁ) — a muscle}

\label{lesson:TE-C251-kandaram}



\begin{quote}
Before the new one: say the Telugu for the mind, then the Telugu for breathlessness.
\end{quote}

\subsection*{You'll want to know: \te{కండరం}}
\textbf{\te{కండరం}} — \emph{kaṇḍaraṁ} — \textquotedblleft{}a muscle\textquotedblright{}.

\begin{grammarlens}[title={What you've built}]
One more word for this chapter.
\end{grammarlens}

\subsection*{Your turn}
\begin{itemize}
\raggedright
  \item \emph{Say it:} \emph{kaṇḍaraṁ}
  \item \emph{Say it:} \emph{kaṇḍaraṁ}, once more
  \item \emph{Say it:} say \emph{kaṇḍaraṁ}
  \item \emph{Recall:} say the Telugu for the mind, then the Telugu for breathlessness, then say \emph{kaṇḍaraṁ} again
\end{itemize}

\begin{tcolorbox}[breakable,colback=teal!4,colframe=teal!35!black,title={Before you move on}]
What does \te{కండరం} mean? (\textquotedblleft{}A muscle\textquotedblright{}.) Say it once more.
\end{tcolorbox}
\section[\texorpdfstring{\te{ఇంజెక్షను} (iñjekṣanu)}{ఇంజెక్షను (iñjekṣanu)}]{\te{ఇంజెక్షను} (iñjekṣanu) — an injection}

\label{lesson:TE-C251-injeksanu}



\begin{quote}
Before the new one: say the Telugu for breathlessness, then the Telugu for a muscle.
\end{quote}

\subsection*{You'll want to know: \te{ఇంజెక్షను}}
\textbf{\te{ఇంజెక్షను}} — \emph{iñjekṣanu} — \textquotedblleft{}an injection\textquotedblright{}.

\begin{grammarlens}[title={What you've built}]
One more word for this chapter.
\end{grammarlens}

\subsection*{Your turn}
\begin{itemize}
\raggedright
  \item \emph{Say it:} \emph{iñjekṣanu}
  \item \emph{Say it:} \emph{iñjekṣanu}, once more
  \item \emph{Say it:} say \emph{iñjekṣanu}
  \item \emph{Recall:} say the Telugu for breathlessness, then the Telugu for a muscle, then say \emph{iñjekṣanu} again
\end{itemize}

\begin{tcolorbox}[breakable,colback=teal!4,colframe=teal!35!black,title={Before you move on}]
What does \te{ఇంజెక్షను} mean? (\textquotedblleft{}An injection\textquotedblright{}.) Say it once more.
\end{tcolorbox}
\section[\texorpdfstring{\te{టీకా} (ṭīkā)}{టీకా (ṭīkā)}]{\te{టీకా} (ṭīkā) — a vaccination}

\label{lesson:TE-C251-tika}



\begin{quote}
Before the new one: say the Telugu for a muscle, then the Telugu for an injection.
\end{quote}

\subsection*{You'll want to know: \te{టీకా}}
\textbf{\te{టీకా}} — \emph{ṭīkā} — \textquotedblleft{}a vaccination\textquotedblright{}.

\begin{grammarlens}[title={What you've built}]
One more word for this chapter.
\end{grammarlens}

\subsection*{Your turn}
\begin{itemize}
\raggedright
  \item \emph{Say it:} \emph{ṭīkā}
  \item \emph{Say it:} \emph{ṭīkā}, once more
  \item \emph{Say it:} say \emph{ṭīkā}
  \item \emph{Recall:} say the Telugu for a muscle, then the Telugu for an injection, then say \emph{ṭīkā} again
\end{itemize}

\begin{tcolorbox}[breakable,colback=teal!4,colframe=teal!35!black,title={Before you move on}]
What does \te{టీకా} mean? (\textquotedblleft{}A vaccination\textquotedblright{}.) Say it once more.
\end{tcolorbox}
